% Contents/lutas_e_combates.tex
\midsectiontitle{Combate em OverGrown}

\section*{A Rolagem de Iniciativa}

No \textit{Overgrown}, o combate não é um mero caos, mas uma sequência de movimentos calculados onde a velocidade e a percepção ditam a sobrevivência. Quando as lâminas se cruzam, a ordem das ações é determinada pelo teste de \link{att:grace}{Graça (Grace)} + Grace Pura. Este valor representa a prontidão do aventureiro em resposta à ameaça iminente.

Aqueles que se destacam agem primeiro, estabelecendo o ritmo do confronto, seguidos pelos demais em ordem decrescente de seus resultados.

\begin{rpgboxtips}{A Estratégia do Atraso}
Um combatente astuto sabe que nem sempre golpear primeiro é a melhor escolha. Você pode optar por \textbf{atrasar sua ação}, aguardando o momento exato em que um aliado abre uma brecha ou um inimigo se expõe. Ao fazer isso, sua posição na ordem de combate é permanentemente redefinida para este novo marco até o desfecho da cena.
\end{rpgboxtips}

\section*{O Fluxo do Tempo e Ações}

O tempo em combate é uma percepção fragmentada. Cada turno condensa entre 0,5 e 6 segundos, onde decisões devem ser tomadas com precisão antes de serem enunciadas. Sob o olhar do mestre, o impossível pode ser tentado, mas apenas o que for coerente com o fluxo temporal será executado.

\subsection*{A Hierarquia das Ações}

A economia de combate em \textit{Overgrown} é baseada na versatilidade. Cada combatente possui uma quantidade limitada de energia que se traduz em três tipos de ações: \textbf{Padrão (ou Completa)}, \textbf{Movimento} e \textbf{Livre}. 

A maestria reside na conversão de esforços. O sistema permite que o jogador sacrifique uma ação de maior complexidade para realizar múltiplas tarefas mais simples.

\begin{itemize}
    \item \textbf{Ação Padrão:} O núcleo do seu turno. Usada para atacar ou usar mágias. Pode ser convertida em uma ação de Movimento, Livre ou reação.
    \item \textbf{Ação de Movimento:} Utilizada para reposicionamento. Pode ser convertida em uma ação Livre.
    \item \textbf{Ação Livre:} Ações rápidas ou frases curtas. Não pode ser convertida.
    \item \textbf{Reação:} Uma resposta instintiva fora do seu turno (1 por rodada). Não pode ser convertida ou trocada.
\end{itemize}

% --- NOVO CONTEÚDO ADICIONADO AQUI ---
\section*{Nível de Ameaça e Precisão}

\originlink{stats:crit}
No mundo de \textit{Overgrown}, o \textbf{Nível de Ameaça} descreve o potencial do aventureiro de encontrar brechas na defesa inimiga e atingir pontos vitais com precisão absoluta. 

Por padrão, um Crítico ocorre ao obter um \textbf{20 natural} nos dados. Entretanto, através de armas específicas e do avanço nos conhecimentos da \textbf{Árvore de Talentos}, este limiar de precisão pode ser reduzido. Um aventureiro de elite pode expandir sua margem de ameaça para resultados como 19, 18 ou até menores, tornando cada golpe uma ameaça potencialmente letal.

\begin{rpgboxtips}{Nós de Agilidade}
Não ignore a Árvore de Talentos durante seus treinos. Existem \textbf{Nós de Estatísticas} e \textbf{Nós Supremos} focados especificamente em quebrar as regras de tempo, permitindo que ações de movimento concedam bônus extras ou que reações sejam recuperadas mais rapidamente.
\end{rpgboxtips}
